\section*{Capítulo \ref{ch:b2:living-soil} --- O solo vivo}

\begin{solution}{exo:b2:living-soil:1}
O: serapilheira e matéria orgânica parcialmente decomposta, onde começa a
decomposição. A: \omterm{def:b2:living-soil:soil}{solo} mineral misturado com \omterm{def:b2:living-soil:soil}{húmus}, escuro e de estrutura
grumosa, que contém a maior parte das raízes, dos organismos, dos
nutrientes e da água. B: o horizonte de acumulação, onde se depositam a
argila, os óxidos de ferro e os carbonatos lixiviados de cima. C:
material de origem intemperizado, fragmentos de rocha numa matriz fina,
onde é liberado o novo \omterm{def:b2:living-soil:soil}{solo} mineral. R: rocha matriz não intemperizada.
\end{solution}

\begin{solution}{exo:b2:living-soil:2}
As cargas negativas das superfícies da argila e do \omterm{def:b2:living-soil:soil}{húmus} retêm cátions
trocáveis que as raízes podem absorver em troca de prótons. Os grãos de
areia são de quartzo, sem carga e grossos, com pouca superfície por
grama e praticamente sem \omterm{def:b2:living-soil:soil}{húmus}. A calagem acrescenta carbonato de
cálcio: o cálcio desloca os íons hidrogênio e alumínio que saturam os
sítios de troca de um \omterm{def:b2:living-soil:soil}{solo} ácido, o carbonato neutraliza a acidez, e os
sítios voltam a ser ocupados por um cátion nutriente.
\end{solution}

\begin{solution}{exo:b2:living-soil:3}
Decompositores: bactérias (\emph{Bacillus}) e fungos
(\emph{Trichoderma}); pastadores microbianos: \omterm{def:b2:unicellular-diversity:protist}{protistas} (amebas) e
nematoides bacterívoros; detritívoros: minhocas, diplópodes, colêmbolos,
tatuzinhos; predadores: ácaros, nematoides predadores, centopeias,
carabídeos; e toupeiras ou musaranhos no topo. Os pastadores, com uma
$\mathrm{C:N}$ mais alta que a de suas presas bacterianas,
excretam o nitrogênio excedente como amônio: eles mineralizam o
nitrogênio que as bactérias tinham imobilizado.
\end{solution}

\begin{solution}{exo:b2:living-soil:4}
Arbuscular: o fungo (um glomeromiceto) penetra nas células do córtex da
raiz e forma arbúsculos; encontrada na maioria das ervas e das culturas
e nas árvores tropicais. Ectomicorriza: o fungo (basidiomicetos e
ascomicetos, muitos deles cogumelos) envolve a ponta da raiz e cresce
entre as células sem penetrar nelas; encontrada nas árvores temperadas e
boreais. Nos dois casos, a planta dá açúcar e o fungo dá fosfato,
nitrogênio, água, micronutrientes e alguma proteção contra patógenos.
\end{solution}

\begin{solution}{exo:b2:living-soil:5}
$L^{*} = I/k = 4/0.8 = \qty{5}{t/ha}$; \omterm{def:b2:biogeochemical-cycles:reservoir}{tempo de residência} $1/k = 1.25$
ano; perda de \qty{90}{\%} após $\ln 10/0.8 = 2.9$ anos.
\end{solution}

\begin{solution}{exo:b2:living-soil:6}
Para cada \qty{100}{kg} de carbono comido, os micróbios retêm
\qty{40}{kg} e precisam de $40/8 = \qty{5}{kg}$ de nitrogênio. A palha
fornece $100/80 =
\qty{1.25}{kg}$: imobilização de \qty{3.75}{kg} por \qty{100}{kg}
de carbono. O trevo fornece $100/15 = \qty{6.7}{kg}$: mineralização
de \qty{1.7}{kg}.
\end{solution}

\begin{solution}{exo:b2:living-soil:7}
\qty{20}{t/ha} de \omterm{def:b2:living-soil:soil}{solo} por ano passam pelas minhocas. Os \qty{20}{cm}
superficiais pesam $0.2\times 10^{4}\times 1.3 = \qty{2600}{t/ha}$:
130 anos para uma passagem. As dez toneladas por acre de Darwin são
\qty{25}{t/ha}, a mesma ordem de grandeza --- seu ``a cada poucos anos''
referia-se à terra vegetal superficial, alguns centímetros, e não a toda
a camada superficial.
\end{solution}

\begin{solution}{exo:b2:living-soil:8}
Bactérias: $10^{9}\times 10^{-12} = \qty{1}{mg}$ por grama. Hifas:
volume $\pi(2\times 10^{-4})^{2}\times 10^{4} = \qty{1.26e-3}{cm^3}$,
massa de \qty{1.4}{mg} por grama. Por hectare ($2.6\times 10^{9}$ g):
\qty{2.6}{t} de bactérias e \qty{3.6}{t} de fungos.
\end{solution}

\begin{solution}{exo:b2:living-soil:9}
$H^{*} = 1.5/0.02 = \qty{75}{t/ha}$ de carbono. Com aração: $1.5/0.04
= \qty{37.5}{t/ha}$; perda de \qty{37.5}{t/ha}, metade dela após
$\ln 2/0.04 = 17$ anos.
\end{solution}

\begin{solution}{exo:b2:living-soil:10}
Fosfato: $\sqrt{2\times 10^{-13}\times 8.64\times 10^{6}} =
\qty{1.3}{mm}$ numa estação --- uma raiz a \qty{1}{cm} da seguinte
alcança um oitavo do \omterm{def:b2:living-soil:soil}{solo} entre elas, enquanto hifas a
\qty{0.1}{mm} umas das outras alcançam todo ele. Nitrato: \qty{4}{cm},
mais que o espaçamento das raízes: ele chega à raiz por si mesmo (e por
fluxo de massa com a água), de modo que as micorrizas acrescentam pouco
para o nitrato e muito para o fosfato.
\end{solution}

\begin{solution}{exo:b2:living-soil:11}
Massa $0.25\times 10^{4}\times 1.3 = \qty{3250}{t/ha}$. Perda líquida de
\qty{19}{t/ha} por ano: tempo de vida de 170 anos. Carbono perdido $20\times
0.03 = \qty{0.6}{t/ha}$ por ano. Sob plantio direto, a perda líquida é de
\qty{1}{t/ha}: 3250 anos.
\end{solution}

\begin{solution}{exo:b2:living-soil:12}
A serapilheira se renova em um ou dois anos, o \omterm{def:b2:living-soil:soil}{húmus} em cinquenta, o
horizonte A mineral em milhares (sua formação é de \qty{0.1}{mm} por ano),
e o sal se acumula em décadas e é drenado em décadas, se houver
drenagem. Um agricultor vê a serapilheira e a cultura responderem em uma
estação e o \omterm{def:b2:living-soil:soil}{húmus} ao longo de uma carreira (um declínio de um terço leva
trinta anos; uma recuperação, o mesmo tanto); a \omterm{prop:b2:living-soil:erosion}{erosão do solo} mineral é
invisível numa vida e irreversível em dez; o sal é visível numa geração
e, sem drenagem, permanente. As variáveis lentas do \omterm{def:b2:living-soil:soil}{solo} são as que
nenhuma contabilidade anual registra.
\end{solution}

\begin{solution}{pb:b2:living-soil:1}
\textbf{1.} $0.25\times 10^{4}\times 1.3 = \qty{3250}{t/ha}$;
carbono $\qty{81}{t/ha}$.
\textbf{2.} $81/12 = \qty{6.8}{t/ha}$ de nitrogênio.
\textbf{3.} Bactérias \qty{1}{mg/g}, $\qty{3.25}{t/ha}$; hifas
$\pi(2\times 10^{-4})^{2}\times 2\times 10^{4}\times 1.1 =
\qty{2.8}{mg/g}$, \qty{9}{t/ha}: biomassa microbiana de cerca de
\qty{12}{t/ha}.
\textbf{4.} Cerca de \qty{14}{t/ha} de matéria viva (massa fresca),
uns \qty{17}{\%} do carbono do \omterm{def:b2:living-soil:soil}{húmus} --- alguns por cento dele em
carbono, pois os organismos são sobretudo água.
\textbf{5.} Carbono microbiano de cerca de metade de \qty{12}{t},
\qty{6}{t}; renovando-se duas vezes por ano com \qty{40}{\%} retidos,
os micróbios consomem $2\times 6/0.4 = \qty{30}{t}$ de carbono e respiram
\qty{18}{t/ha} de carbono, \qty{66}{t} de $\mathrm{CO_2}$ (uma estimativa
grosseira: a retenção e a renovação são aproximadas).
\textbf{6.} Três vezes a produção primária líquida da cultura, de
\qty{5}{t}: a estimativa é alta demais, pois as hipóteses (toda a
biomassa se renovando duas vezes, nada dela dormente) são generosas; a
respiração real do \omterm{def:b2:living-soil:soil}{solo} sob uma cultura é da ordem da própria produção
da cultura, várias toneladas de carbono por hectare por ano. O \omterm{def:b2:living-soil:soil}{solo}
respira tanto quanto as plantas acima dele.
\textbf{7.} Carbono $5\times 0.45 = \qty{2.25}{t}$; nitrogênio $5\times
0.03 = \qty{150}{kg}$; $\mathrm{C:N} = 15$.
\textbf{8.} Mineraliza: os micróbios precisam de $2250\times 0.4/8 =
\qty{112}{kg}$ de nitrogênio e o resíduo contém 150, de modo que
\qty{38}{kg/ha} são liberados em termos líquidos.
\textbf{9.} $1 - e^{-0.5} = \qty{39}{\%}$ após 3 meses; $1 - e^{-2}
= \qty{86}{\%}$ após um ano.
\textbf{10.} $0.39\times 38 = \qty{15}{kg/ha}$.
\textbf{11.} A cultura precisa de \qty{150}{kg}, a maior parte em seus
dois primeiros meses; o resíduo dá \qty{15}{kg} em três meses e
\qty{38}{kg} no total --- um complemento, e não um suprimento. Seu valor
está também no nitrogênio que o trevo fixou e que se encontra nos
próprios \qty{150}{kg} do resíduo, a maior parte dos quais entra no
\omterm{def:b2:living-soil:soil}{húmus} e volta ao longo dos anos.
\textbf{12.} Palha: carbono \qty{2250}{kg}, nitrogênio $2250/80 =
\qty{28}{kg}$; os micróbios precisam de \qty{112}{kg}, de modo que
\qty{84}{kg/ha} são imobilizados a partir do \omterm{def:b2:living-soil:soil}{solo} --- mais da metade da
necessidade da cultura. O agricultor deve acrescentar nitrogênio com a
palha, ou incorporá-la bem antes da semeadura, ou deixá-la na superfície,
onde se decompõe lentamente.
\textbf{13.} Com aração: $1.2/0.025 = \qty{48}{t/ha}$; plantio direto:
$1.2/0.015 = \qty{80}{t/ha}$.
\textbf{14.} $H(t) = 80 - 32\,e^{-0.015t}$: ganho de \qty{4.5}{t/ha}
após 10 anos, \qty{17}{t/ha} após 50.
\textbf{15.} $\mathrm{d}H/\mathrm{d}t = 0.015\times 32 =
\qty{0.48}{t/ha}$ de carbono no primeiro ano, \qty{1.8}{t} de
$\mathrm{CO_2}$.
\textbf{16.} $0.48\times 1.5\times 10^{9} = \qty{0.7}{GtC}$ no
primeiro ano, \qty{7}{\%} das emissões fósseis; no quinquagésimo ano, a
taxa caiu para $0.48\,e^{-0.75} = \qty{0.23}{t/ha}$,
\qty{0.34}{GtC}.
\textbf{17.} O ganho é a aproximação de um novo \omterm{def:b2:biogeochemical-cycles:reservoir}{estado estacionário}: ele
decai até zero à medida que $H$ chega a \qty{80}{t/ha}, e todo o ganho
volta ao ar em décadas se o campo for arado de novo --- um sumidouro no
\omterm{def:b2:living-soil:soil}{solo} é uma transferência única e reversível, e não uma compensação
permanente de emissões contínuas.
\textbf{18.} $k_h = 0.02$: $H^{*} = \qty{60}{t/ha}$, uma perda de
\qty{20}{t/ha}, mais que o ganho dos primeiros cinquenta anos de plantio direto: o aquecimento pode
desfazer o manejo.
\textbf{19.} Com aração: $15 - 0.5 = \qty{14.5}{t/ha}$ por ano; plantio
direto: \qty{1}{t/ha}.
\textbf{20.} $3250/14.5 = 220$ anos; $3250/1 = 3250$ anos.
\textbf{21.} Carbono $15\times 0.025 = \qty{0.38}{t/ha}$; nitrogênio
$0.38/12 = \qty{31}{kg/ha}$ por ano.
\textbf{22.} A decomposição do \omterm{def:b2:living-soil:soil}{húmus} no \omterm{def:b2:biogeochemical-cycles:reservoir}{estado estacionário} sob aração é
$k_hH^{*} =
I = \qty{1.2}{t/ha}$ por ano; a erosão acrescenta mais um terço disso,
mas, ao contrário da decomposição, não é compensada pela entrada ---
ela esgota o estoque.
\textbf{23.} Em parte. O carbono erodido enterrado num sedimento ou num
\omterm{def:b2:biogeochemical-cycles:reservoir}{reservatório} pode ficar protegido da decomposição por muito tempo (um
sumidouro por soterramento); mas os agregados se quebram no transporte,
boa parte do carbono é oxidada no caminho, e a fertilidade perdida do
campo é reposta por fertilizantes cuja fabricação emite carbono. O sinal
global do efeito da erosão sobre o carbono ainda é debatido.
\textbf{24.} Serapilheira $1/k = 0.5$ ano; \omterm{def:b2:living-soil:soil}{húmus} $1/k_h = 40$ anos
sob aração e 67 sob plantio direto; \omterm{def:b2:living-soil:soil}{solo} mineral $3250/0.5 =
6500$ anos em relação à formação, e 220 anos em relação à perda sob
aração.
\textbf{25.} Estoque de carbono de \qty{81}{t/ha} no horizonte A; tempo
de vida de 220 anos sob aração e 3250 sob plantio direto; o resíduo de
trevo fornece cerca de \qty{38}{kg/ha} de nitrogênio mineral, 15 dos
quais nos três primeiros meses.
\end{solution}
